% =====================================================================
% RELATÓRIO DE MACROENTREGA, projetos EMBRAPII
% Modelo para macroentrega intermediária, inclusive a primeira.
% Para o relatório final, use main-final.tex.
%
% Compilar:    make               (ou: latexmk -pdf main.tex)
% Conferir:    make verificar     (log, marcadores, resumo, requisitos)
% Entregar:    make entrega       (marcador pendente vira erro)
%
% Antes de escrever, abra o Anexo I do contrato e copie a descrição
% literal da macroentrega. Ela é o critério de aceite e vai na caixa de
% transcrição da Introdução.
% =====================================================================
\documentclass[rascunho]{macroentrega}

% Logomarcas do cabeçalho, da esquerda para a direita (de uma a quatro).
% O logo do parceiro entra quando o contrato previr.
\setlogos{LOGO/itec_logo.png,LOGO/furg_logo.png,LOGO/embrapii_logo.png}

% ---------------------------------------------------------------------
% DADOS DA CAPA E DA PÁGINA DE ASSINATURAS
% Número e título da macroentrega idênticos aos do SRInfo e do Anexo I.
% Quebre o subtítulo à mão, com linhas de tamanho parecido.
% Campos opcionais: codigo, subtitulo, deliverable, wp, parceiro,
% confidencialidade, codigorelatorio. Apague a linha para ocultar.
% Valor que contém vírgula vai entre chaves: titulo = {A, B}.
% ---------------------------------------------------------------------
\dadosrelatorio{
  tipo              = intermediaria,
  numero            = \ph{NN},
  projeto           = \ph{NOME DO PROJETO},
  codigo            = \ph{número no contrato e no Sisproj},
  subtitulo         = \ph{Primeira linha do subtítulo}\\\ph{segunda linha equilibrada},
  titulo            = \ph{título da macroentrega, como no SRInfo},
  deliverable       = \ph{DN},
  wp                = \ph{WPn},
  parceiro          = \ph{razão social do parceiro tecnológico},
  equipe            = {\ph{Nome Sobrenome},\ph{Nome Sobrenome},
                       \ph{Nome Sobrenome},\ph{Nome Sobrenome}},
  unidade           = {Unidade EMBRAPII iTec/FURG},
  local             = \ph{Cidade},
  data              = \ph{Mês de AAAA},
  % confidencialidade = {CONFIDENCIAL},
  % codigorelatorio   = {SIGLA-RMA-AAAA-NN},
  coordenador       = \ph{Nome do Coordenador},
  razaosocial       = \ph{RAZÃO SOCIAL},
  representante     = \ph{Nome do Representante},
  cargorepresentante= \ph{Cargo},
  dataassinatura    = \ph{DD de mês de AAAA},
}

\begin{document}

\capa
\begin{resumo}

\orientacao{Um parágrafo único, de 150 a 300 palavras, escrito por último, quando o documento já estabilizou. É o único trecho que todo leitor lê e o relatório inteiro comprimido, com resultado e não só com escopo. Diga, nesta ordem: o que a entrega compreende, como foi verificada, o que se obteve (com número e condição de medida) e qual o TRL declarado. Sem superlativo sobre a própria entrega.}

\ph{O que a macroentrega entregou, em uma ou duas frases.} \ph{Como a entrega foi verificada: método, ambiente e quantidade de ensaios.} \ph{Resultado principal, com número e condição de medida.} \ph{Nível de maturidade tecnológica atingido, na forma ``o projeto atinge o nível TRL~X, comprovado por Y''.}

\end{resumo}

\tableofcontents

% Opcionais, úteis a partir de umas cem páginas:
% \listoffigures
% \listoftables
% \begin{siglas}
%   \item[ROS] \textit{Robot Operating System}.
% \end{siglas}

% ---------------------------------------------------------------------
% CAPÍTULOS: um arquivo por capítulo, numerado de dez em dez.
% Para acrescentar uma frente, copie 30_frente_1.tex para 40_frente_2.tex.
% Ao remover um capítulo, não renumere os demais arquivos.
% ---------------------------------------------------------------------
\chapter{Introdução}
\label{chap:introducao}

\orientacao{Uma a duas páginas, nesta ordem: (1) identificar o projeto, o parceiro, o marco contratual e o período; (2) nomear a macroentrega com o texto do Anexo I e dizer o que ela compreende; (3) ligar à entrega anterior, dizendo o que já estava estabelecido e o que esta entrega acrescenta; (4) guiar a leitura, com uma frase por capítulo.

Não abra com contexto de mercado nem com a motivação original do projeto, que pertencem à proposta já aceita. A introdução de uma macroentrega intermediária começa no projeto, não no mundo. Registre em texto que o projeto é realizado com apoio da EMBRAPII, na forma que o contrato exigir.}

O presente relatório documenta a Macroentrega~\ph{NN} do projeto \ph{NOME DO PROJETO}, realizado pela \unidade{} com apoio da EMBRAPII, em parceria com \ph{parceiro tecnológico}, no período de \ph{mês/ano} a \ph{mês/ano}. Esta entrega compreende \ph{descrição da macroentrega, nos termos do Anexo I}.

\ph{Vínculo com a entrega anterior: o que já estava estabelecido na Macroentrega NN-1 e o que esta entrega acrescenta. Obrigatório em macroentrega intermediária.}

O Capítulo~\ref{chap:requisitos} apresenta \ph{o estado dos requisitos e a evolução em relação à entrega anterior}. O Capítulo~\ref{chap:frente1} descreve \ph{o que a frente 1 entregou}. O Capítulo~\ref{chap:discussao} discute \ph{o significado dos resultados e as limitações}. O Capítulo~\ref{chap:conclusao} declara o cumprimento da macroentrega, o nível de maturidade tecnológica atingido e a base deixada para a entrega seguinte.

\section{Objetivos}
\label{sec:objetivos}

\orientacao{Três níveis, sem misturar. O objetivo geral e os objetivos específicos do projeto são transcrição literal do contrato e não mudam de uma entrega para outra: divergência entre relatórios sugere escopo movediço. Os objetivos da macroentrega são o recorte desta entrega dentro deles.

Escreva os objetivos da macroentrega no infinitivo e de modo que o próprio relatório permita verificá-los (``operar o veículo em gramado real'' é verificável; ``avançar na maturidade do sistema'' não é). Se o escopo mudou em relação ao Anexo I, diga aqui, em uma frase, com a referência ao registro da mudança (ata, memorando ou aditivo).

O texto dentro das caixas de transcrição vem do contrato e não se edita, mesmo que a redação pudesse melhorar. Se a fonte não estiver à mão, deixe o marcador e peça o texto à coordenação.}

\subsection{Objetivo Geral do Projeto}

\begin{transcricao}{Anexo I, objetivo geral}
\ph{Transcrever literalmente o objetivo geral do projeto.}
\end{transcricao}

\subsection{Objetivos Específicos do Projeto}

\begin{transcricao}{Anexo I, objetivos específicos}
\begin{enumerate}[label=\alph*)]
    \item \ph{Transcrever literalmente o primeiro objetivo específico};
    \item \ph{Transcrever literalmente o segundo objetivo específico}.
\end{enumerate}
\end{transcricao}

\subsection{Descrição da Macroentrega~\numeromacroentrega{} no Anexo I}

\begin{transcricao}{Anexo I, macroentrega \numeromacroentrega}
\ph{Transcrever literalmente a descrição contratual da macroentrega, igual ao texto do SRInfo.}
\end{transcricao}

\subsection{Objetivos da Macroentrega~\numeromacroentrega}

\ph{Frase de abertura: o que esta macroentrega se propôs a entregar, em redação própria.} Os objetivos são:
\begin{enumerate}[label=\alph*)]
    \item \ph{verbo no infinitivo e objeto verificável pelo próprio relatório};
    \item \ph{verbo no infinitivo e objeto verificável pelo próprio relatório}.
\end{enumerate}

\chapter{Requisitos do Projeto}
\label{chap:requisitos}

\orientacao{Capítulo obrigatório na primeira macroentrega, que estabelece a lista completa, e atualizado em todas as seguintes com uma coluna a mais. Derive os requisitos dos objetivos e da metodologia do plano de trabalho, separe funcionais de não funcionais e dê um identificador único a cada um.

A coluna \textbf{Verificação} impede que a tabela vire declaração. Nas linhas \parc{} e \done{}, informe o método (ensaio, inspeção, análise ou demonstração) e o ponteiro para a evidência no relatório (Tab.~4.2, Fig.~5.1, Anexo~A). \done{} significa concluído e verificado. Requisito derivado ou alterado depois da primeira entrega entra na tabela com o registro da mudança, e a tabela nunca some entre uma entrega e outra.

Na primeira macroentrega, use uma única coluna, \texttt{\{M01\}}, e troque a seção de evolução por um parágrafo de leitura da linha de base. A lista de requisitos da primeira macroentrega é o instrumento que mede o projeto inteiro.

Para ajustar as colunas, edite o argumento de \texttt{tabelarequisitos}, por exemplo \texttt{\{M01,M02,M03\}}. Cada linha usa \texttt{\textbackslash req\{ID\}\{descrição\}\{status por coluna\}\{verificação\}}.}

Este capítulo apresenta os requisitos do projeto \ph{NOME DO PROJETO} atualizados para a Macroentrega~\numeromacroentrega, que definem as capacidades, as restrições e os atributos de qualidade que o sistema deve satisfazer ao longo do desenvolvimento. Os requisitos foram derivados dos objetivos e da metodologia do plano de trabalho e estão organizados em funcionais, que descrevem o que o sistema deve ser capaz de fazer, e não funcionais, que definem como ele deve operar. Cada requisito tem identificador único, descrição, status em cada macroentrega e a verificação que sustenta o status.

O status é indicado pelas seguintes siglas: \todo{} (não iniciado), \wip{} (em andamento), \parc{} (parcialmente implementado) e \done{} (concluído e verificado).

\section{Requisitos Funcionais}
\label{sec:req_funcionais}

\subsection{\ph{Área de competência} (RF-XXX)}

\begin{tabelarequisitos}{M01,M02}
\req{RF-XXX-01}{\ph{O sistema deve ser capaz de ..., em condição observável.}}{\todo & \wip}{}
\req{RF-XXX-02}{\ph{O sistema deve ..., com valor e unidade quando houver meta numérica.}}{\wip & \done}{Ensaio, Fig.~\ref{fig:exemplo}}
\req{RF-XXX-03}{\ph{O sistema deve ..., verificável por inspeção ou análise.}}{\parc & \parc}{Análise, Tab.~\ref{tab:exemplo}}
\end{tabelarequisitos}

\section{Requisitos Não Funcionais}
\label{sec:req_nao_funcionais}

\subsection{\ph{Atributo de qualidade} (RNF-XXX)}

\begin{tabelarequisitos}{M01,M02}
\req{RNF-XXX-01}{\ph{O sistema deve atender ..., com limite e condição de medida.}}{\todo & \todo}{}
\req{RNF-XXX-02}{\ph{O software deve ser versionado em repositório com documentação de interfaces.}}{\wip & \wip}{}
\end{tabelarequisitos}

\section{Evolução dos Requisitos}
\label{sec:evolucao_requisitos}

\orientacao{Leitura em prosa da tabela, com contagens verificáveis na própria tabela, nunca estimativas: quantos requisitos mudaram de status, onde o avanço se concentrou, o que permanece \todo{} e por qual dependência. O comando \texttt{make verificar} imprime a contagem por status e por coluna. Na primeira macroentrega, esta seção vira um parágrafo de leitura da linha de base.}

Dos \ph{N} requisitos definidos para o projeto, \ph{N} tiveram o status alterado entre a Macroentrega~\ph{NN-1} e a Macroentrega~\numeromacroentrega, e \ph{N} estão concluídos e verificados. O avanço concentrou-se em \ph{áreas ou identificadores}. Os requisitos ainda não iniciados vinculam-se a \ph{área}, cuja execução depende de \ph{dependência técnica}.

\chapter{\ph{Nome da frente ou do eixo}}
\label{chap:frente1}

\orientacao{Um capítulo por frente de trabalho, ou por eixo quando a entrega atravessa frentes. Copie este arquivo para cada frente nova e mantenha a mesma estrutura em todas: reduz a variação entre capítulos escritos por pessoas diferentes.

Regras para o capítulo técnico:
\begin{itemize}
    \item Descreva o que existe, não o que se pretendia. Passado para o trabalho realizado, presente para o sistema como ele é.
    \item Decisão de projeto vem com o motivo (``optou-se por X porque Y''). Comparação de alternativas é conteúdo legítimo.
    \item Toda afirmação de desempenho vem com a condição de medida: quantos ensaios, em que cena, com que objeto, sob que parâmetros.
    \item Resultado negativo se escreve como resultado, com a causa. Falha é informação técnica de valor.
    \item Apresente o dado antes de interpretá-lo, para que o leitor possa discordar da interpretação sem descartar o dado.
    \item Componente usado por mais de uma frente é descrito uma vez e referenciado depois.
\end{itemize}

Na primeira macroentrega, um dos capítulos técnicos traz a prova de conceito que sustenta o TRL declarado: o que foi montado ou implementado, em que condição foi testado e o que se observou. Prova de conceito em simulação serve, desde que traga figura, dado e condição experimental declarada.}

\section{Introdução}
\label{sec:frente1_intro}

\ph{O que esta frente entrega nesta macroentrega e qual objetivo da macroentrega ela atende.}

\section{Objetivos da Frente nesta Entrega}

\ph{Objetivos da frente, no infinitivo, recortados dos objetivos da macroentrega.}

\section{Metodologia e Desenvolvimento}
\label{sec:frente1_metodologia}

\ph{O que foi feito e como, no passado; o sistema como ele é, no presente. Cada decisão de projeto com o motivo.}

\section{Setup Experimental}
\label{sec:frente1_setup}

\ph{Ambiente, equipamentos, parâmetros e procedimento, em detalhe suficiente para reproduzir o ensaio.} Os parâmetros constam na Tabela~\ref{tab:parametros}.

\begin{table}[htbp]
\centering
\caption{\ph{Parâmetros do ensaio, com unidade.}}
\label{tab:parametros}
\estilotabela
\begin{tabularx}{\linewidth}{|P{4cm}|C{3cm}|X|}
\hline
\rowcolor{reqcolor}\cab{Parâmetro} & \cab{Valor} & \cab{Observação} \\ \hline
\ph{parâmetro} & \ph{valor e unidade} & \ph{condição} \\ \hline
\rowcolor{reqlight}\ph{parâmetro} & \ph{valor e unidade} & \ph{condição} \\ \hline
\end{tabularx}
\end{table}

\section{Resultados}
\label{sec:frente1_resultados}

\orientacao{Evidência forte: figura do artefato operando no ambiente descrito, tabela de resultados com condição experimental declarada, gráfico de erro ou de trajetória, registro de execução com data e local, parâmetros que permitam reproduzir o experimento. Diagrama de arquitetura e captura de tela mostram intenção ou existência, não funcionamento.

Cada objetivo da macroentrega precisa apontar para ao menos uma evidência forte. Se só houver evidência fraca, o texto diz que o objetivo foi cumprido em parte.

Figuras: vetoriais (PDF) sempre que possível, com rótulos legíveis no tamanho impresso, sem moldura nem título de slide. Toda figura e toda tabela é chamada no texto antes de aparecer, e a chamada diz o que observar. A legenda diz o que a figura mostra e não repete o título da seção.}

A Figura~\ref{fig:exemplo} apresenta \ph{o que a figura mostra e o que o leitor deve observar nela}.

\begin{figure}[htbp]
\centering
% Figura real: \includegraphics[width=\linewidth]{IMAGES/nome_da_figura.pdf}
\figuraph[5cm]{figura}
\caption{\ph{Legenda que diz o que a figura mostra, com a condição do ensaio.}}
\label{fig:exemplo}
\end{figure}

A Tabela~\ref{tab:exemplo} resume \ph{os resultados e a condição em que foram obtidos}.

\begin{table}[htbp]
\centering
\caption{\ph{Resultados por condição de ensaio.}}
\label{tab:exemplo}
\estilotabela
\begin{tabularx}{\linewidth}{|P{3.6cm}|C{2.6cm}|C{2.6cm}|X|}
\hline
\rowcolor{reqcolor}\cab{Condição} & \cab{Ensaios} & \cab{Resultado} & \cab{Observação} \\ \hline
\ph{condição} & \ph{n} & \ph{valor, unidade} & \ph{nota} \\ \hline
\rowcolor{reqlight}\ph{condição} & \ph{n} & \ph{valor, unidade} & \ph{nota} \\ \hline
\end{tabularx}
\end{table}

\section{Discussão}
\label{sec:frente1_discussao}

\ph{O que o resultado desta frente significa, sob quais condições vale e o que ele não demonstra.}

% \include{capitulos/40_frente_2}
% \include{capitulos/50_resultados_consolidados}
\chapter{Discussão e Limitações}
\label{chap:discussao}

\orientacao{Capítulo curto e obrigatório. Responde a três perguntas: o que o resultado significa, sob quais condições vale e o que ainda não está demonstrado.

Escreva cada limitação com a causa técnica e o efeito sobre o resultado, em parágrafo próprio com título em negrito. Registre o fato, sem culpado e sem pedido de desculpa: ``o fornecimento do hardware ocorreu em junho, o que concentrou os ensaios em duas semanas'' é registro; ``a demora do fornecedor inviabilizou o trabalho'' atribui culpa.

Limitação declarada com sobriedade aumenta a credibilidade do documento e delimita por escrito o que não foi prometido. Limitação omitida reaparece na reunião de aceite, quando admiti-la custa muito mais.}

\section{Discussão}

\ph{Interpretação integrada dos resultados dos capítulos técnicos, depois de os dados terem sido apresentados.}

\section{Limitações}

\paragraph{\ph{Nome da limitação}.} \ph{Causa técnica da limitação e efeito sobre o resultado.}

\paragraph{\ph{Nome da limitação}.} \ph{Causa técnica da limitação e efeito sobre o resultado.}

\chapter{Considerações Finais}
\label{chap:conclusao}

\orientacao{Meia a uma página de texto, mais os quadros. Quatro funções: declarar o cumprimento retomando a descrição do Anexo I; sintetizar o que foi alcançado em cada frente, uma frase por frente; declarar o nível de maturidade no início e no fim, dizendo qual dimensão mudou e apontando a evidência; indicar a base deixada para a entrega seguinte.

Não traga informação nova, não repita a discussão e não termine com frase de efeito sobre o futuro da tecnologia.}

\section{Cumprimento da Macroentrega}
\label{sec:cumprimento}

A Macroentrega~\numeromacroentrega, referente a ``\ph{descrição da macroentrega, igual à do Anexo I}'', foi \ph{concluída / concluída em parte, com a parte faltante identificada}. A Tabela~\ref{tab:cumprimento} relaciona cada cláusula da descrição contratual ao ponto do relatório que a responde e à evidência que a sustenta.

\orientacao{Quebre a descrição do Anexo I em cláusulas verificáveis, uma por linha. Cláusula sem resposta é lacuna de entrega e não se resolve com redação: leve à coordenação antes de escrever. Resultado sem cláusula correspondente é resultado adicional: descreva-o depois de demonstrar o escopo contratado e identifique-o como adicional, sem que ocupe o lugar de item contratado.}

\begin{table}[htbp]
\centering
\caption{Cumprimento da descrição da macroentrega no Anexo I.}
\label{tab:cumprimento}
\estilotabela
\begin{tabularx}{\linewidth}{|P{4.0cm}|P{2.4cm}|X|C{1.9cm}|}
\hline
\rowcolor{reqcolor}\cab{Cláusula do Anexo I} & \cab{Resposta} & \cab{Evidência} & \cab{Situação} \\ \hline
\ph{cláusula 1} & Seção~\ref{sec:frente1_resultados} & Fig.~\ref{fig:exemplo} & \done \\ \hline
\rowcolor{reqlight}\ph{cláusula 2} & \ph{Seção} & \ph{Fig. ou Tab.} & \parc \\ \hline
\end{tabularx}
\end{table}

\section{Síntese por Frente}

\begin{itemize}
    \item \ph{Frente 1: uma frase com o artefato ou achado e o número que o caracteriza.}
    \item \ph{Frente 2: uma frase com o artefato ou achado e o número que o caracteriza.}
\end{itemize}

\section{Maturidade Tecnológica}
\label{sec:trl}

\orientacao{Escreva na forma ``o projeto atinge o nível X, comprovado por Y'', e não ``o projeto atingiu o nível X''. O valor declarado precisa ficar dentro da faixa prevista no Anexo I e ser o mesmo lançado no SRInfo.

A evidência cobre as três dimensões da norma ABNT NBR ISO 16290: escala do objeto, fidelidade do objeto e ambiente em que a função tecnológica é desenvolvida. Defina o ambiente antes de usar o rótulo: uma frase com o envelope de condições que o Anexo I exige (iluminação, terreno, velocidade, volume de dados, perfil de uso) e outra com o subconjunto em que o ensaio desta macroentrega ocorreu. A distância entre as duas é o que o nível declarado precisa suportar.

Declare o nível no início e no fim e diga qual dimensão mudou. Se nenhuma mudou, provavelmente não houve avanço de nível, e declarar TRL mantido é resultado legítimo. Em entrega de integração, vale o nível do componente menos maduro do caminho crítico, e convém nomeá-lo. Ensaio em bancada não sustenta ``ambiente relevante''.}

O projeto atinge o nível TRL~\ph{X}, comprovado por \ph{evidência}. Os níveis de maturidade são caracterizados por escala do objeto, fidelidade do objeto e ambiente, conforme a norma ABNT NBR ISO 16290~\cite{abnt16290}. O ambiente exigido pelo Anexo I é \ph{envelope de condições}; o ensaio desta macroentrega ocorreu em \ph{subconjunto dessas condições}.

\begin{table}[htbp]
\centering
\caption{Nível de maturidade tecnológica no início e no final da macroentrega.}
\label{tab:trl}
\estilotabela
\begin{tabularx}{\linewidth}{|P{2.6cm}|X|X|P{2.6cm}|}
\hline
\rowcolor{reqcolor}\cab{Dimensão} & \cab{Início (TRL~\ph{X})} & \cab{Final (TRL~\ph{Y})} & \cab{Evidência} \\ \hline
Escala do objeto & \ph{escala} & \ph{escala} & \ph{Fig. ou Tab.} \\ \hline
\rowcolor{reqlight}Fidelidade do objeto & \ph{fidelidade} & \ph{fidelidade} & \ph{Fig. ou Tab.} \\ \hline
Ambiente & \ph{ambiente} & \ph{ambiente} & \ph{Fig. ou Tab.} \\ \hline
\end{tabularx}
\end{table}

\section{Próximos Passos}

\orientacao{Direção técnica em duas ou três frentes, cada uma com a dependência que a condiciona. Sem cronograma, sem responsável e sem métrica: tarefa, dono e prazo pertencem à ata e à proposta da entrega seguinte, e escrevê-los aqui cria compromisso não contratado. Prefira ``o passo seguinte é ampliar a coleta'' a ``o sistema passará a generalizar para novas situações''.}

\ph{Base estabelecida por esta entrega e direção técnica da entrega seguinte, em verbo modesto.}


\referencias{referencias}

\appendix
\chapter{\ph{Título do anexo}}
\label{anx:a}

\orientacao{Anexo é o lugar do material que interrompe a leitura mas precisa ficar registrado: tabelas longas de parâmetros e de resultados brutos, lista de componentes com especificação, endereços de repositórios com a versão entregue (\emph{commit} ou \emph{tag}) e registros de propriedade intelectual. Não cole trecho extenso de código; aponte o repositório. Cada anexo é citado no corpo do relatório. Para tabela que atravessa páginas, use \texttt{longtable}, que repete o cabeçalho.}

\ph{Conteúdo do anexo, citado no corpo do relatório.}


\paginaassinaturas

\end{document}
